

\medskip
\noindent\SGARefTarget{sga3:iv:proposition:4:4:9}\textbf{Proposition 4.4.9.}
\textit{Every equivalence relation in \(\widetilde{\mathscr C}\) is
universally effective (\SGARefLink{sga3:iv:definition:3:3:3}{3.3.3}). More precisely, let \(R\) be a
\(\widetilde{\mathscr C}\)-equivalence relation on the sheaf \(X\). Then the
sheaf associated with the separated presheaf}
\[
  i(X)/i(R):S\longmapsto X(S)/R(S)
\]
\textit{is a universal effective quotient of \(X\) by \(R\).}

Let \(X/R\) denote the sheaf quotient of \(X\) by \(R\), which exists by
\SGARefLink{sga3:iv:theorem:4:4:1:item:ii:01}{4.4.1(ii)}:
\[
  X/R=a\bigl(i(X)/i(R)\bigr).
\]
We must show that \(X\to X/R\) is a universal effective epimorphism and that
the morphism
\[
  f:R\longrightarrow X\times_{X/R}X
\]
is an isomorphism. The first assertion has already been proved in \SGARefLink{sga3:iv:proposition:4:4:3}{4.4.3}.
As for \(f\), it is obtained by applying the functor \(a\) to the morphism
\[
  i(R)\longrightarrow i(X)\times_{i(X/R)}i(X),
\]
or, since \(i(X)/i(R)\) is separated by \SGARefLink{sga3:iv:lemma:4:4:5:item:ii:01}{4.4.5(ii)}, so that
\[
  i(X)/i(R)\longrightarrow i(X/R)
\]
is a monomorphism, to the canonical morphism
\[
  i(R)\longrightarrow
  i(X)\times_{\,i(X)/i(R)}i(X).
\]
Thus we are reduced to proving the same statement in the category of
presheaves. But \(i(X)/i(R)\) is the presheaf
\[
  S\longmapsto X(S)/R(S),
\]
so \SGARefLink{sga3:iv:proposition:4:4:9}{the assertion} reduces to its set-theoretic analogue, where it is
immediate.

\medskip
\noindent\SGARefTarget{sga3:iv:proposition:4:4:10}\textbf{Proposition 4.4.10.}
\textit{Under the hypotheses of \SGARefLink{sga3:iv:proposition:4:4:9}{4.4.9}, let \(Y\) be a subsheaf of \(X\), and
write \(R_Y\) for the equivalence relation induced on \(Y\) by \(R\). Then the
canonical morphism (\SGARefLink{sga3:iv:definition:3:1:6}{3.1.6})}
\[
  Y/R_Y\longrightarrow X/R
\]
\textit{is a monomorphism. It identifies \(Y/R_Y\) with a subsheaf of
\(X/R\), namely the sheaf image of the composite}
\[
  Y\longrightarrow X\longrightarrow X/R.
\]

The morphism of presheaves
\[
  i(Y)/i(R_Y)
  =i(Y)/i(R)_{i(Y)}
  \longrightarrow i(X)/i(R)
\]
is a monomorphism. Since the functor \(a\) is left exact (\SGARefLink{sga3:iv:remark:4:3:16}{4.3.16}), it carries
monomorphisms to monomorphisms; hence
\[
  Y/R_Y\longrightarrow X/R
\]
is a monomorphism. The final assertion follows from the commutative diagram
\[
\begin{array}{ccc}
  Y&\longrightarrow&X\\
  \big\downarrow&&\big\downarrow\\
  Y/R_Y&\longrightarrow&X/R
\end{array}
\]
and the fact that \(Y\to Y/R_Y\) is covering.

In view of \SGARefLink{sga3:iv:proposition:4:4:10}{this proposition}, we shall always identify \(Y/R_Y\) with a
subsheaf of \(X/R\).

